\begin{statement}{C.2.UN.78}{Proved}{Exact root, birational action and coarsening records for the Hopf-minimum branch}
Let $f:X\to\Delta_t$ be a proper degeneration of connected
complex two-tori from a smooth threefold, with arbitrary
central support $S=f^{-1}(0)_{\rm red}$. Take an actual
embedded resolution $\pi:Y\to X$ by smooth central
point and curve blowups, and write
\[
 \operatorname{div}(f\pi)=\sum_i m_iD_i,
 \qquad \operatorname{div}(\pi^*\Omega)=\sum_i k_iD_i,
 \qquad r_i=\frac{k_i+1}{m_i},
 \qquad N=\operatorname{lcm}_i(m_i),
\tag{C.RK.1}
\]
where $\Omega$ generates $f_*K_X$. The $D_i$ are the
smooth components of the resolved SNC support. Form
the actual normalization and a toric semistable
resolution
\[
 \begin{gathered}
 Y_N=\operatorname{Norm}(Y\times_{\Delta_t}\Delta_s),
                \qquad t=s^N,\\
 Z\xrightarrow{\ q\ }Y_N\xrightarrow{\ \rho\ }Y,
                 \qquad G=\mu_N.
 \end{gathered}
\tag{C.RK.2}
\]
These maps can be chosen $G$-equivariantly, with $Z$
smooth and full reduced simple-NC central fibre.
All components added by $q$ are ruled or rational.

The following is an explicit finite-cover condition
on this resolution record: every strict root-cover
component whose smooth minimal model is Hopf lies
over an index $i$ with
\[
                            r_i=\min_j r_j.
\tag{C.RK.3}
\]
For exactly the records satisfying (C.RK.3), the
semistable root $Z$ admits the finite compatible
canonical-flat reduction of UN75. Denote its
actual comparison by $\phi:Z\dashrightarrow X_0$.
The flat endpoint has one of the three complete
core types in UN75, and it carries the actual
finite \emph{birational} action
\[
                 \gamma_0=\phi\gamma\phi^{-1}
                 \quad (\gamma\in G).
\tag{C.RK.4}
\]
Its group laws hold as equalities of meromorphic
maps; it acts over $s\mapsto\zeta s$ and is regular
on the punctured family. Retain the full record
$\phi$, the toric coarsening $q$, and $\pi$.
Their inverse and quotient operations reconstruct
the given original degeneration exactly:
\[
 X_0\ \dashrightarrow\ Z\ \xrightarrow{q}\ Y_N
       \ \longrightarrow\ Y_N/G=Y
       \ \xrightarrow{\pi}\ X.
\tag{C.RK.5}
\]
In particular the inverse record is applied
\emph{before} the finite quotient. No quotient of
an undefined central birational action is used.
Conversely the same operations on any actual
flat-core record with these retained regularization,
root-coarsening and smooth blowdown maps produce
a torus degeneration in the stated branch.
This gives complete finite generation for this
explicit resolution-and-root branch, without
an equivariant canonical-flat reduction assertion.

The Euler correction for the root is explicit.
For a nonempty index set $I$ with nonempty stratum
\[
 D_I^\circ=
    \left(\bigcap_{i\in I}D_i\right)
           \setminus\bigcup_{j\notin I}D_j,
 \qquad d_I=\gcd_{i\in I}m_i,
 \qquad \nu_I=\frac{N^{|I|-1}d_I}{\prod_{i\in I}m_i},
\tag{C.RK.6}
\]
one has $1\le|I|\le3$ and $\nu_I$ is a positive
integer. Put
\[
 \begin{split}
 R_N&=\sum_I(d_I\nu_I-1)e_c(D_I^\circ),\\
 \Delta_\pi&=2p_\pi+
                   \sum_\alpha(2-2g(C_\alpha)),
 \end{split}
\tag{C.RK.7}
\]
where $p_\pi$ and $C_\alpha$ are the actual point
and connected curve centres of $\pi$. Let
$\Delta_K$ be the signed normal-cone defect of
the forward UN75 record from $Z$ to $X_0$.
Then the exact original-support formula is
\[
                e(S)=e(S_0)-\Delta_K-R_N-\Delta_\pi.
\tag{C.RK.8}
\]
Thus zero Euler number is equivalent to
$\Delta_K+R_N+\Delta_\pi=e(S_0)$, where the
right side is zero for the smooth or marked
cycle core, and $2[N_0:\Gamma_0]$ for the
toric core. The symbols $N_0,\Gamma_0$ here
are its two-dimensional lattices, separate
from the root exponent $N$ in (C.RK.1).
Formula (C.RK.8) agrees with the universal
original-support formula of UN39 by the
actual root and blowup maps.

The test (C.RK.3) is in fact intrinsic to
the original degeneration. Write its original
full fibre and generator orders as
$\sum_\alpha n_\alpha S_\alpha$ and
$\sum_\alpha a_\alpha S_\alpha$.
On the normalization $T_\alpha$ of $S_\alpha$
form the finite normal cover $C_\alpha$
by the following actual local equations
and transition maps:
\[
 \begin{gathered}
 f=g_j^{\,n_\alpha}u_j,\qquad
 w_j^{\,n_\alpha}=u_j|_{T_\alpha},\\
 w_k=(g_j/g_k)|_{T_\alpha}\,w_j .
 \end{gathered}
\tag{C.RK.18}
\]
Here $g_j$ is a local defining equation
of the Cartier component $S_\alpha$;
the polynomial covers are normalized
after gluing. Define
\[
 \tau=\sup\left\{c\in\R:
       |f|^{-2c}|\Omega|^2\text{ is locally integrable
                              near }S\right\}.
\tag{C.RK.19}
\]
The notation $|\Omega|^2$ denotes its
positive top-degree density, so this
definition requires no metric choice.
Then (C.RK.3) is equivalent to
\[
 \frac{a_\alpha+1}{n_\alpha}=\tau
 \quad\text{for every }\alpha\text{ for which a connected
 component of }C_\alpha\text{ has Hopf minimal model}.
\tag{C.RK.20}
\]
Consequently its truth is independent
of the embedded resolution, its
semistable triangulation and the
chosen root exponent divisible by
all resolved multiplicities.

For an original compact family the record can
be made on a finite compact cover of its base,
and (C.RK.5) recovers the original compact
family. Reverse generation on compact actual
core records with the displayed finite action
and coarsening likewise remains compact.
The original endpoint neighbourhood, not only
its central log space, is retained throughout.
\end{statement}

\begin{proof}
\par\medskip\noindent\textbf{Embedded resolution and the exact normalized root charts.}
Use analytic embedded resolution by central
smooth centres, as in UN39 and
\cite[Theorems~1.6 and1.10]{BierstoneMilman1997}.
Divisor centres give identity blowups and are
omitted. The remaining centres in a smooth
threefold are points and compact curves.
The actual pullback of $\Omega$ is the
canonical generator on $Y$: apply the
Hartogs direct-image comparison (C.KG.6)
to the ordinary blowup maps.

Near a stratum with $r$ incident components,
absorb the unit in the fibre equation into
a coordinate. The exact local equation and
volume form have the shape
\[
 t=x_1^{m_1}\cdots x_r^{m_r},\qquad
 \pi^*\Omega=u\,x_1^{k_1}\cdots x_r^{k_r}
                         dx_1\wedge dx_2\wedge dx_3,
                         \quad u\text{ a unit}.
\tag{C.RK.9}
\]
The normalization of $t=s^N$ has $d_I$
local branches. Each is the affine toric
chart with cocharacter lattice
\[
 \mathcal L_I=\left\{(v,h)\in\Z^r\oplus\Z:
                         \sum_i m_iv_i=Nh\right\}
\tag{C.RK.10}
\]
and cone generated by the height-one rays
$((N/m_i)e_i,1)$. The omitted $3-r$
coordinates are unchanged. Local roots of
the unit, if not absorbed in (C.RK.9),
give the same normalization maps and
the actual finite transition characters.
The group $G$ permutes the $d_I$ branches
and acts torically on each stabilised
branch, over $s\mapsto\zeta s$.

The height-one slice has dimension at most
two. Subdivide its edges at every lattice
point, and triangulate each two-dimensional
lattice triangle using all its lattice
points. An empty lattice triangle is
unimodular, by Pick's formula. Thus these
subdivisions give smooth toric charts.
They agree on their subdivided edges.
The refinement can be obtained by
successively inserting the height-one
lattice points as star subdivisions.
Each step is the normalized blowup
of its actual monomial ideal. On an
overlap, replacing a boundary equation
by a holomorphic unit times that
equation multiplies every monomial
generator by a unit and leaves this
ideal unchanged. Thus the normalized
blowup maps glue through the original
tangent-dependent transition units;
those transitions have not been
replaced by constant torus maps.
The stabilizer of an old coarse cone fixes
each old ray: the original SNC components
are individually labelled, and the group
acts over their identity map in $Y$.
Its action on this lattice slice is
therefore the identity. Choose the
triangulations orbitwise. The actual
toroidal transition maps, which preserve
these labelled rays, glue the subdivisions
to a proper $G$-equivariant map $q$.

Every new ray still has height one,
so $s$ has order one on every central
component of $Z$. The resulting toric
boundary is simple NC with smooth
components. For the embedded assertion
retain every old component and stratum
label as well as each subdivision ray.
Transition units fix those labels and
the ray vector; distinct rays in one
old cone are never identified. The
local components in a unimodular
chart therefore belong to distinct
global components, and their smooth
star charts glue without a global
self-incidence. Over a double curve the
added surfaces are projective-line
bundles over finite covers of that
curve, with their indicated modifications;
over a triple point they are rational
toric surfaces. None has Hopf minimal
model. The normalized fibre product
over $\Delta_s^*$ is connected because
its compact torus fibres and its base
are connected. The proper Stein factor
has degree one there and hence on the
normal disc. Thus $Z\to\Delta_s$ has
connected full central fibre as well.

\par\medskip\noindent\textbf{Canonical ratios and the explicit minimum test.}
At a primitive height-one ray $(v,1)$,
the toric Jacobian of (C.RK.9) gives
\[
 \operatorname{ord}_{(v,1)}(q^*\rho^*\pi^*\Omega)
                 =\sum_i(k_i+1)v_i-1.
\tag{C.RK.11}
\]
For a strict ray this is
$N(k_i+1)/m_i-1$. Any height-one
point is a convex combination of the
old vertices, so (C.RK.11) is the
same affine combination of their
orders. The minimum among all rays is
therefore
\[
                  a_{\min}=N\min_i r_i-1.
\tag{C.RK.12}
\]
It is an integer, since a minimum is
attained at an old vertex. The form
\[
            \Omega_Z=s^{-a_{\min}}
                             q^*\rho^*\pi^*\Omega
\tag{C.RK.13}
\]
is holomorphic and generates the
canonical direct image on the reduced
semistable family. Its component
orders have minimum zero. Since all
new components are ruled or rational,
UN75's Hopf-minimum condition is
exactly (C.RK.3) for the strict covers.
This proves both directions of the
test for the chosen resolution record.
It also retains the actual generator
character
$\gamma^*\Omega_Z=\zeta^{-a_{\min}}\Omega_Z$.

\par\medskip\noindent\textbf{The actual finite birational action and its quotient order.}
Apply UN75 to $Z$. Its comparison
$\phi$ is a bimeromorphic map over
$\Delta_s$, with a specified inverse
record and an actual punctured
isomorphism. Formula (C.RK.4)
therefore defines meromorphic maps
on $X_0$. Their products satisfy
the group laws on the dense
punctured set and consequently
as meromorphic maps. They have
order $N$ on the total family,
as seen on the base. Choosing
the endpoint volume by
$\phi^*\Omega_0=\Omega_Z$ also
gives exactly the character in
the preceding paragraph under
this birational action.

To reconstruct, first reverse
the retained $\phi$-record.
The action is then the original
regular $G$-action on $Z$.
Apply its equivariant toric
coarsening $q$ to get the actual
normal space $Y_N$. Its quotient
is literally $Y$. Indeed on a
normal local ring $A$ of $Y$,
the root ring is the integral
closure $B$ of
$A[s]/(s^N-f\pi)$. The invariants
of its total ring of fractions are
$\operatorname{Frac}(A)$.
An invariant element of $B$ is
integral over $A$ and lies in
that fraction field, hence lies
in $A$. Thus
\[
                              B^G=A.
\tag{C.RK.14}
\]
This argument applies to the
normalization on the connected
total space; the local split
branches are permuted by $G$
and give the same invariant
ring equality. The retained
ordinary blowdowns $\pi$ now
recover $X$. These exact maps
prove (C.RK.5) and its converse.
An equivariant replacement of
$\phi$ has not been used.

\par\medskip\noindent\textbf{The root correction from actual stratum fibres.}
Over $D_I^\circ$ the reduced
central support of $Y_N$ is an
\emph{actual} finite etale cover
of degree $d_I$. The local root
equation in (C.RK.9) proves the
degree; unit-root transition
maps glue the cover and may
permute its sheets. At a point
of this cover the toric fibre
of $q$ has Euler number $\nu_I$.
To see this, use its torus-orbit
stratification. All positive-
dimensional torus orbits have
Euler number zero. The zero-
dimensional orbits are exactly
the maximal unimodular cones
in the subdivision, so their
number is the normalized
volume of the height-one
simplex.

Here is its exact lattice
volume. Set $a_i=m_i/d_I$ and
choose an integral vector $b$
with $\sum_i a_i b_i=1$.
The difference lattice of
the slice is
$\ker(a:\Z^r\to\Z)$.
Appending $b$ to a basis of
this kernel gives a unimodular
basis of $\Z^r$. Consequently
the normalized volume is the
absolute determinant obtained
by appending $b$ to the $r-1$
vertex differences. Expansion
gives
\[
 \left|\det\left[
 \frac{N}{m_1}e_1-\frac{N}{m_r}e_r,
 \ldots,
 \frac{N}{m_{r-1}}e_{r-1}-\frac{N}{m_r}e_r,
 b\right]\right|
               =\frac{N^{r-1}d_I}{\prod_i m_i}.
\tag{C.RK.15}
\]
For $r=1$ both sides are one.
This proves the formula and
integrality of $\nu_I$ without
a choice of triangulation.

Proper stratified Euler
pushforward now gives
\[
 \begin{split}
 e(S_Z)&=\sum_I d_I\nu_I e_c(D_I^\circ),\\
 e(S_Y)&=\sum_I e_c(D_I^\circ),\\
 e(S_Z)-e(S_Y)&=R_N.
 \end{split}
\tag{C.RK.16}
\]
This also checks the finite
quotient directly. A deck
element $\gamma^j$ fixes the
$d_I$-sheeted central cover
over this stratum exactly
when $d_I\mid j$. Therefore
\[
 \begin{split}
 e(S_{Y_N}^{\gamma^j})
       &=\sum_{d_I\mid j}d_I e_c(D_I^\circ),\\
 \frac1N\sum_{j=0}^{N-1}
          e(S_{Y_N}^{\gamma^j})
       &=\sum_I e_c(D_I^\circ)=e(S_Y).
 \end{split}
\tag{C.RK.17}
\]
All these equalities hold for
the actual covers even when
they are globally connected
and their sheet labels have
nontrivial monodromy.

The original resolution has
$e(S_Y)-e(S)=\Delta_\pi$ by
its ordinary point and curve
maps in UN39. The forward
UN75 record has
$e(S_0)-e(S_Z)=\Delta_K$.
Combine these two equalities
with (C.RK.16) to obtain
(C.RK.8), including its zero
criterion. Each summand has
been compared through an
actual proper map or a
retained overlap-compatible
record.

\par\medskip\noindent\textbf{Intrinsic covers and the intrinsic minimum.}
The component $S_\alpha$ is Cartier
in the original smooth threefold.
Its local equations satisfy
$g_k=v_{kj}g_j$ for holomorphic
units, and therefore
$u_k=v_{kj}^{-n_\alpha}u_j$.
This is precisely the transition
in (C.RK.18), including its cocycle
identity. The corresponding monic
finite algebras glue before taking
their finite normalizations.
On the open smooth component
away from other components this
is the $\mu_{n_\alpha}$-torsor
of roots of the normal-line
trivialization supplied by $f$.
No global meromorphic frame for
that normal line is chosen.

For any root exponent divisible
by $n_\alpha$, the exact map from
the strict central root cover is
\[
                         w_j=\frac{s^{N/n_\alpha}}{g_j}.
\tag{C.RK.21}
\]
At its generic point numerator
and denominator have the same
order, and their quotient satisfies
(C.RK.18). On a smooth strict
component of $Z$ this meromorphic
quotient is integral over its
holomorphic local ring, by that
monic equation. Normality makes
it holomorphic. The map therefore
extends to a proper map to the
glued finite cover, and factors
through its normalization. It
has generic degree one onto a
connected component of $C_\alpha$.
Every such component is obtained,
since on the dense open both
covers have the same complete
set of $n_\alpha$ roots.
Thus the strict surfaces are
actual birational resolutions
of these fixed covers for every
choice of embedded resolution
and divisible root exponent.

Every exceptional surface of
the original embedded resolution
is projective: initially it is
$\PP^2$ or a projective-line
bundle over a compact curve,
and its subsequent transforms
are projective modifications.
The ideals on these projective
surfaces are algebraic by GAGA.
A finite normal analytic cover
of a projective surface is also
projective: its finite coherent
algebra and multiplication
algebraize by GAGA, and the
finite pullback of an ample
line is ample. Its surface
resolutions remain bimeromorphic
to projective surfaces. Such
root components cannot have
Hopf minimal model. The new
semistable components have
already been shown ruled or
rational. Hence (C.RK.21)
accounts for every Hopf root
component, with no exceptional
ones omitted.

Finally apply the literal change
of variables under $\pi$ to the
density in (C.RK.19), outside
the exceptional sets of measure
zero. Near a resolved crossing
it is a bounded positive unit
times the product of the radial
densities
$|x_i|^{2(k_i-cm_i)}\,dA(x_i)$,
with smooth tangent density.
On finitely many relatively
compact charts about the compact
central fibre, the integral is
finite exactly when
$k_i-cm_i>-1$ for every $i$.
At equality it diverges already
at a general point of that
divisor. Thus
\[
                             \tau=\min_i\frac{k_i+1}{m_i}.
\tag{C.RK.22}
\]
This is the original density's
integrability threshold, so it
is independent of resolution.
The strict transform of an
original component retains
its original pair
$(a_\alpha,n_\alpha)$.
Together with (C.RK.21) and the
exceptional-component comparison,
(C.RK.22) proves exactly
(C.RK.20) and the stated
choice independence.

\par\medskip\noindent\textbf{Compact scope.}
For a compact base $B$ choose
a meromorphic function with
a simple zero at the selected
point; Riemann--Roch supplies
one with its poles away from
that point. Normalize the
cyclic cover $s^N=h$ of $B$.
It is a compact curve with
the exact local parameter
map used in (C.RK.2).
Additional branch values
are retained in this base
record. Perform the normal
base change and equivariant
resolution globally, and
the UN75 replacement near
the chosen central fibre.
Its actual annular map
glues to the unchanged
outside as in UN75. Conjugate
the global finite action
through that comparison;
it is a finite birational
action on this compact
family. Reversing the
record regularizes it,
and the compact proper
coarsening and quotient
in (C.RK.5) recover the
original compact family.
The converse uses those
same proper maps on an
actual compact core record.
No replacement of a full
neighbourhood follows only
from the central-log-space
classification of UN71
or its occurrence theorem
UN73.
\end{proof}


% C.2.UN.74 from part_C2_ruled_crossing_contraction.tex

